% Part: lambda-calculus
% Chapter: introduction
% Section: basic-pr-lambda

\documentclass[../../../include/open-logic-section]{subfiles}

\begin{document}

\olfileid{lam}{int}{bas}
\olsection{ప్రాథమిక ఆదిమ పునరావృత్త ప్రమేయాలు \usetoken{S}{lambda definable}}

\begin{lem}
$\Zero$, $\Succ$, $\Proj{n}{i}$ ప్రమేయాలు \tetoken{లాంబ్డాతో నిర్వచించదగిన}{!!{lambda definable}}వి.
\end{lem}

\begin{proof}
$\Zero$ అంటే కేవలం
$\lambd[x][\lambd[y][y]]$.

ఉత్తరగామి ప్రమేయం~$\Succ$ను
$\fn{Succ}(u) = \lambd[x][\lambd[y][x(uxy)]]$ ద్వారా నిర్వచిస్తాం.
ఇది ఎందుకు పనిచేస్తుందో ఆలోచించాలి: పునరావర్తకంగా భావించిన ప్రతి
సంఖ్యాంకం~$\num{n}$కు, ప్రతి ప్రమేయం~$f$కు,
$\fn{Succ}(\num{n},f)$ అనేది ఇన్‌పుట్~$y$పై $y$తో మొదలుపెట్టి $f$ను
$n$ సార్లు ప్రయోగించి, ఆ తరువాత మరోసారి ప్రయోగించే ప్రమేయం.

ప్రక్షేపాల గురించి చెప్పడానికి ఏమీ లేదు: $\fn{Proj}^n_i(x_0, \dots,
x_{n-1}) = x_i$. మరో మాటలో, మన సంకేత సంప్రదాయాల ప్రకారం
$\fn{Proj}^n_i$ అనేది $\lambd[x_0][\dots \lambd[x_{n-1}][x_i]]$
అనే లాంబ్డా పదం.
\end{proof}

\end{document}
